% !TEX TS-program = XeLaTeX
% use the following command:
% all document files must be coded in UTF-8
\documentclass[portuguese]{textolivre}
% build HTML with: make4ht -e build.lua -c textolivre.cfg -x -u article "fn-in,svg,pic-align"

\usepackage{float}
\journalname{Texto Livre}
\thevolume{19}
%\thenumber{1} % old template
\theyear{2026}
\receiveddate{\DTMdisplaydate{2025}{9}{9}{-1}} % YYYY MM DD
\accepteddate{\DTMdisplaydate{2025}{10}{24}{-1}}
\publisheddate{\DTMdisplaydate{2026}{4}{2}{-1}}
\corrauthor{Marcelo Magalhães Foohs}
\articledoi{10.1590/1983-3652.2026.62236}
%\articleid{NNNN} % if the article ID is not the last 5 numbers of its DOI, provide it using \articleid{} commmand 
% list of available sesscions in the journal: articles, dossier, reports, essays, reviews, interviews, editorial
\articlesessionname{articles}
\runningauthor{Foohs, Isele e Rebouças} 
%\editorname{Leonardo Araújo} % old template
\sectioneditorname{Daniervelin Pereira~\orcid{0000-0003-1861-3609}}
\layouteditorname{Leonardo Araújo~\orcid{0000-0003-3884-2177}}

\title{Criatividade na formação docente: narrativas digitais interativas e pensamento computacional}
\othertitle{Creativity in teacher education: interactive digital narratives and computational thinking}
% if there is a third language title, add here:
%\othertitle{Artikelvorlage zur Einreichung beim Texto Livre Journal}

\author[1]{Marcelo Magalhães Foohs~\orcid{0000-0002-4735-0732}\thanks{Email: \href{mailto:mmfoohs@gmail.com}{mmfoohs@gmail.com}}}
\author[1]{Priscila Castioni Isele~\orcid{0000-0002-1040-6608}\thanks{Email: \href{mailto:priscila.castioni@hotmail.com}{priscila.castioni@hotmail.com}}}
\author[1]{Randerson Oliveira Melville Rebouças~\orcid{0009-0005-3056-5074}\thanks{Email: \href{randerson.melville@gmail.com}{randerson.melville@gmail.com}}}
\affil[1]{Universidade Federal do Rio Grande do Sul, CINTED/PPGIE, Porto Alegre, RS, Brasil.}

\addbibresource{article.bib}
% use biber instead of bibtex
% $ biber article

% used to create dummy text for the template file
\definecolor{dark-gray}{gray}{0.35} % color used to display dummy texts
\usepackage{lipsum}
\SetLipsumParListSurrounders{\colorlet{oldcolor}{.}\color{dark-gray}}{\color{oldcolor}}

% used here only to provide the XeLaTeX and BibTeX logos
\usepackage{hologo}

% if you use multirows in a table, include the multirow package
\usepackage{multirow}

% provides sidewaysfigure environment
\usepackage{rotating}

% CUSTOM EPIGRAPH - BEGIN 
%%% https://tex.stackexchange.com/questions/193178/specific-epigraph-style
\usepackage{epigraph}
\renewcommand\textflush{flushright}
\makeatletter
\newlength\epitextskip
\pretocmd{\@epitext}{\em}{}{}
\apptocmd{\@epitext}{\em}{}{}
\patchcmd{\epigraph}{\@epitext{#1}\\}{\@epitext{#1}\\[\epitextskip]}{}{}
\makeatother
\setlength\epigraphrule{0pt}
\setlength\epitextskip{0.5ex}
\setlength\epigraphwidth{.7\textwidth}
% CUSTOM EPIGRAPH - END

% LANGUAGE - BEGIN
% ARABIC
% for languages that use special fonts, you must provide the typeface that will be used
% \setotherlanguage{arabic}
% \newfontfamily\arabicfont[Script=Arabic]{Amiri}
% \newfontfamily\arabicfontsf[Script=Arabic]{Amiri}
% \newfontfamily\arabicfonttt[Script=Arabic]{Amiri}
%
% in the article, to add arabic text use: \textlang{arabic}{ ... }
%
% RUSSIAN
% for russian text we also need to define fonts with support for Cyrillic script
% \usepackage{fontspec}
% \setotherlanguage{russian}
% \newfontfamily\cyrillicfont{Times New Roman}
% \newfontfamily\cyrillicfontsf{Times New Roman}[Script=Cyrillic]
% \newfontfamily\cyrillicfonttt{Times New Roman}[Script=Cyrillic]
%
% in the text use \begin{russian} ... \end{russian}
% LANGUAGE - END

% EMOJIS - BEGIN
% to use emoticons in your manuscript
% https://stackoverflow.com/questions/190145/how-to-insert-emoticons-in-latex/57076064
% using font Symbola, which has full support
% the font may be downloaded at:
% https://dn-works.com/ufas/
% add to preamble:
% \newfontfamily\Symbola{Symbola}
% in the text use:
% {\Symbola }
% EMOJIS - END

% LABEL REFERENCE TO DESCRIPTIVE LIST - BEGIN
% reference itens in a descriptive list using their labels instead of numbers
% insert the code below in the preambule:
%\makeatletter
%\let\orgdescriptionlabel\descriptionlabel
%\renewcommand*{\descriptionlabel}[1]{%
%  \let\orglabel\label
%  \let\label\@gobble
%  \phantomsection
%  \edef\@currentlabel{#1\unskip}%
%  \let\label\orglabel
%  \orgdescriptionlabel{#1}%
%}
%\makeatother
%
% in your document, use as illustraded here:
%\begin{description}
%  \item[first\label{itm1}] this is only an example;
%  % ...  add more items
%\end{description}
% LABEL REFERENCE TO DESCRIPTIVE LIST - END


% add line numbers for submission
%\usepackage{lineno}
%\linenumbers

\begin{document}
\maketitle

\begin{polyabstract}
\begin{abstract}
O artigo investiga manifestações da criatividade na produção de narrativas digitais interativas desenvolvidas por professores da educação básica em uma formação continuada que articula o ensino de Língua Portuguesa e o desenvolvimento do Pensamento Computacional, por meio do Ciclo de Remidiação Ativa (CRA). A abordagem metodológica é qualitativa, de natureza descritiva e interpretativa, tendo como corpus narrativas autorais produzidas no Scratch a partir da reescrita criativa de fábulas. A análise fundamenta-se em aportes da criatividade \cite{torrance1995}, dos letramentos múltiplos \cite{rojo2012} e da aprendizagem criativa mediada \cite{sawyer2004}. Os resultados evidenciam estratégias autorais que reconfiguram expressivamente os textos originais e indicam integração entre linguagem, abstração algorítmica e programação. Conclui-se que o CRA favorece processos criativos e interdisciplinares na formação docente, embora os dados se limitem às produções entregues.

\keywords{Criatividade docente \sep Narrativas digitais interativas \sep Pensamento computacional \sep Letramentos digitais \sep Formação de professores}
\end{abstract}

\begin{english}
\begin{abstract}
This article investigates manifestations of creativity in the production of interactive digital narratives developed by basic education teachers during a continuing education program that articulates Portuguese language teaching and the development of Computational Thinking through the Active Remediation Cycle (ARC). A qualitative, descriptive, and interpretative methodological approach was adopted, with the analytical corpus consisting of authorial narratives created in Scratch based on the creative rewriting of fables. The analysis is grounded in theoretical contributions on creativity \cite{torrance1995}, multiliteracies \cite{rojo2012}, and mediated creative learning \cite{sawyer2004}. The results reveal authorial strategies that expressively reconfigure the original texts and indicate an integration of language, algorithmic abstraction, and programming. It is concluded that the ARC fosters creative and interdisciplinary processes in teacher education, although the data are limited to the submitted productions.

\keywords{Teacher creativity \sep Interactive digital narratives \sep Computational thinking \sep Digital literacies \sep Teacher education}
\end{abstract}
\end{english}
% if there is another abstract, insert it here using the same scheme
\end{polyabstract}

\section{Introdução}\label{sec-intro}

A formação de professores frente aos desafios da cultura digital exige propostas pedagógicas que integrem linguagem, autoria, tecnologia e criatividade de forma articulada. Nesse cenário, narrativas digitais interativas vêm se consolidando como recursos promissores para promover experiências autorais que mobilizam múltiplas linguagens: verbal, algorítmica e de programação visual. A Base Nacional Comum Curricular \cite{bncc2018}, sua versão complementar de Computação \cite{cne2022} e a Nota Técnica 21 do \textcite{cieb2024} reforçam essa direção ao destacar a importância da autoria textual, do pensamento computacional e da criticidade em contextos mediados por tecnologias.

Com base nesses marcos, foi implementada, no primeiro semestre de 2025, uma proposta formativa em uma rede municipal de ensino, estruturada em nove módulos ao longo de três meses. O percurso incluiu encontros presenciais e atividades assíncronas mediadas pelo Google Sala de Aula, organizando-se em torno do Ciclo de Remidiação Ativa (CRA): um processo que articula reescrita de fábulas, decomposição narrativa e programação em Scratch, integrando multiletramentos e pensamento computacional em uma abordagem autoral e processual.

Este artigo analisa processos criativos emergentes nas diferentes linguagens mobilizadas durante a formação. A investigação centra-se na manifestação da criatividade nas produções dos professores em formação, com base em uma matriz que articula os critérios de \textcite{torrance1995} (fluência, originalidade, elaboração e flexibilidade), os indicadores de criatividade computacional propostos por \textcite{brennan2012} e as diretrizes da Nota Técnica 21 do \textcite{cieb2024}.

Essa abordagem permite uma leitura sensível à complexidade das narrativas digitais interativas, compreendidas como formas singulares de autoria mediadas por programação. O objetivo geral é analisar a criatividade nas produções autorais resultantes do CRA. Como objetivos específicos, busca-se:

\begin{enumerate}[label=\alph*)]
\item Identificar e categorizar elementos criativos presentes nas fábulas e narrativas digitais, com base nos critérios de \textcite{torrance1995} e nos indicadores de criatividade computacional de \textcite{brennan2012}.
\item Relacionar tais elementos aos objetivos de aprendizagem da BNCC \cite{bncc2018}, especialmente nas habilidades EF69LP15, EF69LP16 e EF69LP17 (Língua Portuguesa), e da BNCC Complementar de Computação \cite{cne2022}, nos códigos EF69CO02, EF69CO03 e EF69CO05.
\item Refletir sobre o potencial das narrativas digitais interativas como estratégia para o desenvolvimento da criatividade docente e computacional, com vistas à sua aplicação no ensino fundamental.
\end{enumerate}

Por fim, ressalta-se que a proposta foi concebida com intencionalidade replicativa: espera-se que, ao vivenciar o CRA, os professores estejam aptos a adaptá-lo às suas realidades escolares, promovendo com seus alunos práticas criativas e autorais alinhadas às competências da BNCC \cite{bncc2018,cne2022}.

\section{Fundamentação teórica}\label{sec-fundamentacao}

Para compreender os fundamentos que orientam esta pesquisa, é necessário abordar diferentes eixos conceituais que sustentam a proposta do Ciclo de Remidiação Ativa (CRA) e o foco analítico na criatividade docente. O percurso teórico parte da concepção de narrativas digitais interativas como expressão multimodal e autoral; passa pela mediação pedagógica na perspectiva histórico-cultural de \textcite{vygotsky2007}; e articula os conceitos de remidiação e multiletramentos como base metodológica do CRA. Na sequência, apresenta-se o CRA como proposta formativa e o Scratch como ambiente de autoria digital, seguido pela discussão dos referenciais utilizados na análise da criatividade: os critérios de \textcite{torrance1995} e os indicadores computacionais de \textcite{brennan2012}.

\subsection{Narrativas digitais interativas na educação}

Narrativas digitais interativas são formas contemporâneas de contar histórias que combinam linguagem verbal, imagens, sons e elementos de programação para criar experiências não lineares e participativas. Diferenciam-se das narrativas digitais tradicionais por ampliarem o papel do leitor, que deixa de ser apenas receptor passivo e assume o papel de usuário, capaz de interagir com a trama, modificar caminhos, realizar escolhas e até interferir no desfecho da história \cite{ryan2009}.

Segundo \textcite{jenkins2006}, a interatividade em ambientes digitais redefine a relação entre autor, texto e leitor, criando ecossistemas narrativos que favorecem o engajamento, a experimentação e a produção cultural distribuída. Nessa perspectiva, as narrativas digitais interativas deslocam o lugar do leitor para uma posição mais ativa e demandam novas formas de letramento.

\textcite{miller2004} reforça essa ideia ao afirmar que tais narrativas exigem do narrador competências multimodais, pois mobilizam linguagens distintas e recursos técnicos que expandem suas possibilidades expressivas. A autora define narrativas digitais interativas como ``[...] materiais narrativos que chegam ao seu público por meio de tecnologia e mídia digitais. Uma de suas marcas únicas é a interatividade — comunicações de mão dupla entre o público e o material narrativo''\footnote{“[\ldots] narrative material that reaches its audience via digital technology and media. One of its unique hallmarks is interactivity — back and forth communications between the audience and the narrative material” \cite[p.~4]{miller2004}.} \cite[p.~4]{miller2004}. Essa combinação entre texto, imagem, som e ação programada torna essas experiências narrativas multimodais e, portanto, pedagogicamente promissoras, sobretudo quando integradas a práticas autorais na escola.

No contexto educacional, narrativas digitais interativas vêm sendo exploradas como estratégia didática que estimula o pensamento crítico, a criatividade e o protagonismo discente. Ao permitir que os estudantes construam mundos narrativos interativos, com personagens, regras e percursos alternativos, esse gênero contribui para o desenvolvimento de competências relacionadas à linguagem, à lógica e à resolução de problemas, dialogando diretamente com os fundamentos dos multiletramentos e do pensamento computacional.

Neste estudo, o conceito de narrativa digital interativa está diretamente associado à produção autoral de professores em formação, realizada no ambiente Scratch. Ao compor histórias com elementos visuais, verbais e programáveis, os docentes exploraram diferentes modos de significação e reconfiguraram fábulas em projetos interativos mediados por algoritmos. Assim, essas narrativas constituem objetos culturais complexos, nos quais a criatividade se manifesta por meio da linguagem, da lógica e do design digital.

\subsection{Mediação e desenvolvimento na perspectiva histórico-cultural}

A perspectiva histórico-cultural do desenvolvimento humano, proposta por \textcite{vygotsky2007}, oferece uma base epistemológica sólida para compreender a aprendizagem como processo social mediado por artefatos e signos. Nessa abordagem, o desenvolvimento cognitivo ocorre por meio da internalização de práticas sociais culturalmente organizadas, mediadas pela linguagem e por artefatos simbólicos \cite{vygotsky2007}, o que inclui, nas sociedades contemporâneas, dispositivos digitais.

Trata-se, portanto, de um conceito central na teoria vygotskiana. O professor, nessa visão, é agente de mediação que organiza o ambiente, promove interações e sustenta processos de significação. A aprendizagem acontece na chamada zona de desenvolvimento proximal (ZDP), entendida como o campo de possibilidades em que o indivíduo, por si só, ainda não é capaz de realizar determinadas tarefas, mas pode desenvolvê-las com a mediação de um outro mais experiente ou de artefatos culturais e tecnológicos. Nessa zona, a mediação pedagógica cumpre papel estratégico na ampliação das possibilidades cognitivas, afetivas e culturais dos sujeitos.

Essa concepção dialoga com \textcite{sawyer2004}, que entende a criatividade como prática sociocultural, emergente das interações entre sujeitos, artefatos e contextos simbólicos. Ao diferenciar ``ensinar com criatividade'' de ``ensinar para a criatividade'', o autor destaca o papel do professor como mediador da expressão autoral e da experimentação estética, especialmente em ambientes abertos e culturalmente significativos. Como ele afirma: ``Ensino criativo é uma improvisação disciplinada, pois sempre ocorre dentro de estruturas e enquadramentos amplos''\footnote{“Creative teaching is disciplined improvisation because it always occurs within broad structures and frameworks” \cite[p.~13]{sawyer2004}).} \cite[p.~13]{sawyer2004}). Nesse sentido, a criatividade docente não se limita à espontaneidade, mas se concretiza por meio de mediações estruturadas, em diálogo com os contextos culturais dos aprendizes, o que converge com a perspectiva vygotskiana e com a dinâmica formativa do CRA.

No contexto do Ciclo de Remidiação Ativa (CRA), a mediação assume papel indispensável em todas as etapas da formação, desde a reescrita textual até a programação digital. A proposta demanda do professor-formador a capacidade de provocar deslocamentos interpretativos, promover reflexão metalinguística e apoiar a transposição de ideias para múltiplas linguagens. Ao mesmo tempo, exige acompanhamento próximo das dificuldades enfrentadas pelos docentes em formação, que lidam simultaneamente com os desafios da produção textual, do pensamento computacional e do uso criativo das tecnologias.

Nesse sentido, artefatos digitais como o Scratch devem ser compreendidos, à luz de \textcite{vygotsky2007}, como artefatos cognitivos, mediadores simbólicos que reconfiguram as práticas de autoria e aprendizagem. O ambiente Scratch permite tanto a implementação de algoritmos e narrativas quanto a reorganização de conceitos, a articulação entre forma e conteúdo e a experimentação criativa em linguagem natural, visual e computacional. Ele se insere, portanto, na tradição vygotskiana de artefatos culturais que expandem os modos de pensamento e expressão dos sujeitos.

Assim, o CRA se ancora na concepção de que o desenvolvimento da criatividade, da linguagem e da lógica não ocorre de forma isolada ou espontânea, mas sim por meio da mediação intencional, contínua e situada. Essa mediação é profundamente cultural, pois opera com signos, linguagens e ferramentas que pertencem a práticas sociais diversas e que, quando apropriadas de forma consciente e reflexiva, ampliam a capacidade de ação simbólica dos professores em formação.

\subsection{Criatividade e pensamento computacional}

A criatividade, entendida como a capacidade de gerar ideias novas, resolver problemas de forma original e expressar-se de maneira singular, constitui uma competência transversal fundamental na formação docente contemporânea. No campo educacional, a criatividade não deve ser concebida como um talento inato ou restrito às artes, mas como uma habilidade mobilizável, ensinável e avaliável em diferentes contextos de aprendizagem. Quando articulada ao uso de tecnologias digitais e à lógica computacional, a criatividade manifesta-se em formas híbridas, que combinam linguagem, estrutura, interação e design.

No contexto do Ciclo de Remidiação Ativa (CRA), a criatividade é compreendida como um processo autoral que se desenvolve ao longo de sucessivas remidiações, nas quais ideias são reformuladas, explicitadas e reorganizadas à medida que transitam entre diferentes formas de expressão. Esse percurso pode iniciar-se por distintos pontos, como a socialização de uma ideia, a análise de um protótipo ou a escrita inicial de um texto, mas se constitui formativamente na passagem consciente entre linguagem natural escrita, decomposição algorítmica e implementação computacional. É nesse movimento que a criatividade se articula de modo indissociável ao Pensamento Computacional (PC), entendido não como habilidade genérica de raciocínio, mas como competência específica de formalização e organização de sentidos.

Embora amplamente difundido no campo educacional, o conceito de Pensamento Computacional ainda é marcado por ambiguidades conceituais. Definições excessivamente amplas tendem a reduzi-lo a qualquer forma de raciocínio estruturado ou resolução de problemas, o que dilui sua identidade computacional. Neste estudo, adota-se uma compreensão do PC alinhada às diretrizes da BNCC Computação \cite{cne2022} e à Nota Técnica nº 21 do \textcite{cieb2024}, segundo a qual o Pensamento Computacional envolve a formulação de soluções por meio de representações estruturadas e formalizáveis, com potencial de serem compreendidas e executadas por máquinas, ainda que expressas em linguagens intermediárias acessíveis ao contexto escolar.

Essa definição permite preservar a especificidade técnica do Pensamento Computacional sem comprometer sua aplicabilidade pedagógica. No CRA, o PC não é tratado como um fim em si mesmo, mas como um meio estruturante que orienta a transformação criativa das narrativas. Ao converter ideias ainda implícitas em textos coesos, algoritmos narrativos ou protótipos interativos, os participantes são levados a explicitar relações lógicas, selecionar informações relevantes e reorganizar estruturas narrativas, exercitando abstração, decomposição, reconhecimento de padrões e construção de algoritmos.

Nessa direção, \textcite{kuz2023}, ao analisar produções desenvolvidas no Scratch a partir de uma abordagem de programação estruturada, reforça a compreensão do Pensamento Computacional como uma competência vinculada à formalização lógica, à organização interna das soluções e à possibilidade de execução computacional. Nessa perspectiva, o Scratch deixa de ser compreendido apenas como um ambiente lúdico ou expressivo e passa a configurar-se como um meio privilegiado para tornar visíveis processos de abstração, decomposição, controle de fluxo, paralelismo e depuração. Essa leitura dialoga diretamente com os indicadores propostos por \textcite{brennan2012} e converge com o Ciclo de Remidiação Ativa (CRA), no qual a transposição consciente entre linguagem natural escrita, estrutura algorítmica e implementação computacional opera como gesto metodológico central. Assim, o CRA não dilui o Pensamento Computacional em práticas criativas difusas, mas o mobiliza de forma situada, autoral e tecnicamente consistente, articulando rigor estrutural e expressividade narrativa.

Estudos recentes indicam que o Pensamento Computacional pode ser mobilizado no ensino de Língua Portuguesa mesmo em contextos não digitais. \textcite{sassi2023}, ao relatarem uma experiência com computação desplugada no Ensino Fundamental, evidenciam que atividades linguísticas envolvendo ordenação, decomposição e reconstrução textual favorecem habilidades como sequenciação lógica, abstração conceitual e reconhecimento de padrões, especialmente como estratégia introdutória para iniciantes ou para crianças que ainda não operam no nível de abstração lógica algorítmica. As autoras ancoram sua análise na definição de \textcite{wing2006}, preservada com variações em textos subsequentes, segundo a qual ``O pensamento computacional se fundamenta no poder e nos limites dos processos computacionais, sejam eles executados por um ser humano ou por uma máquina''\footnote{“Computational thinking builds on the power and limits of computing processes, whether they are executed by a human or by a machine” \cite[p.~33]{wing2006}.} \cite[p.~33]{wing2006}. O ponto desestruturante dessa formulação reside no deslocamento do critério algorítmico: ao admitir como computacionais processos executados exclusivamente por humanos que não carregam, em sua essência, condição semântica de execução por máquina, a definição dissocia o algoritmo de sua própria condição computacional. Não está em jogo a forma de expressão, algoritmos podem ser descritos em linguagem natural ou pseudocódigo, mas a exigência de que contenham, em sua essência semântica, a possibilidade de compreensão e execução por uma máquina.

Assim, a aceitação acrítica da definição de \textcite{wing2006}, sem problematizar o operador lógico OR presente na formulação original, amplia excessivamente o escopo do Pensamento Computacional e dilui sua especificidade algorítmica. Ao admitir como computacionais processos que não possuem condições para serem compreendidos ou executados por sistemas computacionais, corre-se o risco de reduzi-lo a estratégias cognitivas humanas genéricas, desvinculadas de critérios explícitos de executabilidade e controle formal. Nesse sentido, embora a experiência relatada por \textcite{sassi2023} evidencie a potência pedagógica de práticas linguísticas estruturadas, ela também exemplifica uma armadilha recorrente do campo: a naturalização de processos cognitivos exclusivamente humanos como computacionais sem explicitar os critérios que os tornariam passíveis de execução por uma máquina.

No contexto do Ciclo de Remidiação Ativa (CRA), o trabalho linguístico, inclusive em modalidades desplugadas, é compreendido como etapa formativa legítima, mas não como substituto da remidiação computacional propriamente dita. O CRA assume que a mobilização plena do Pensamento Computacional se intensifica quando a linguagem é tensionada em direção à formalização explícita, à decomposição algorítmica e à possibilidade de execução computacional, preservando a distinção entre pensamento humano em geral e seu subconjunto denominado Pensamento Computacional enquanto prática algorítmica situada, autoral e tecnicamente consistente. Portanto, o CRA pressupõe sujeitos que já possuam a capacidade de abstração lógica algorítmica.

No que se refere à análise das produções computacionais, este estudo dialoga com os indicadores propostos por \textcite{brennan2012}, que organizam o Pensamento Computacional no Scratch em três dimensões inter-relacionadas: (a) práticas computacionais, como iterar, modularizar e depurar; (b) conceitos computacionais, como eventos, paralelismo, variáveis e operadores; e (c) perspectivas computacionais, relacionadas à autoria e à expressividade, incluindo a produção para si e a produção orientada ao outro. No contexto do CRA, essas dimensões não se manifestam de forma isolada, mas se articulam progressivamente à medida que os participantes passam da exploração autorreferente para a construção de narrativas computacionais comunicáveis, mobilizando revisão, depuração e reescrita como exigências de endereçamento ao outro.

A decomposição narrativa constitui um momento privilegiado de transição entre criatividade literária e pensamento algorítmico, exigindo que os participantes transformem eventos narrativos em sequências lógicas, condições e relações de causalidade. De modo complementar, a reconstrução de textos em linguagem natural a partir de protótipos ou algoritmos já implementados evidencia que o Pensamento Computacional pode ser mobilizado também em movimento reverso, contribuindo para o desenvolvimento da escrita, da coesão textual e da consciência sobre o uso das variedades linguísticas, inclusive da norma padrão \cite{faraco2015}.

Nesse sentido, a criatividade computacional observada nas produções analisadas não se reduz à originalidade estética ou ao uso expressivo de recursos visuais, mas envolve decisões autorais sustentadas por estruturas lógicas coerentes e representações computacionais consistentes. A criatividade manifesta-se, portanto, na capacidade de articular linguagem, narrativa e computação de modo integrado, equilibrando expressividade e rigor estrutural.

A perspectiva sociocultural de \textcite{sawyer2004} contribui para compreender esse processo como prática emergente, situada e mediada culturalmente. No CRA, a criatividade não é um atributo individual isolado, mas resulta da interação entre sujeitos, artefatos digitais, repertórios culturais e mediações pedagógicas. O Pensamento Computacional atua, nesse contexto, como mediador cognitivo que possibilita a transformação progressiva de ideias em produções autorais multimodais, sem que a dimensão técnica se sobreponha à expressiva.

Assim, ao integrar criatividade e Pensamento Computacional no Ciclo de Remidiação Ativa, este estudo evidencia que a programação em blocos pode funcionar como espaço legítimo de autoria narrativa e desenvolvimento linguístico, desde que ancorada em critérios conceituais claros e práticas pedagógicas que valorizem expressividade e formalização. Essa articulação reforça o potencial do CRA como metodologia formativa voltada ao enfrentamento de fragilidades no uso da língua portuguesa e à promoção do Pensamento Computacional como competência curricular.


\subsection{O Ciclo de Remidiação Ativa (CRA): linguagem e computação articuladas}

O Ciclo de Remidiação Ativa (CRA) constitui uma metodologia ativa desenvolvida a partir de estudos realizados no âmbito de um projeto de pós-doutorado, cujo foco foi investigar formas de articular, de maneira sistemática, o desenvolvimento da linguagem natural escrita e do Pensamento Computacional por meio da autoria de narrativas digitais com o uso do Scratch \cite{foohs2022}. Sua formulação emerge de uma Revisão Integrativa de Literatura e de entrevistas exploratórias com docentes da Educação Básica, que evidenciaram lacunas recorrentes nas práticas pedagógicas envolvendo narrativas digitais, especialmente no que diz respeito à ausência de uma etapa robusta de produção textual prévia e à fragmentação entre escrita e programação.

Estudos recentes têm explorado a articulação entre programação em blocos e contação de histórias como estratégia para promover o pensamento computacional, especialmente no contexto do Scratch \cite{santos2024}. Em muitos desses trabalhos, contudo, a narrativa é construída predominantemente de forma oral ou informal, sendo transposta diretamente para o ambiente computacional sem uma etapa sistemática de produção textual escrita em norma padrão. Embora tais abordagens apresentem resultados positivos em termos de motivação, criatividade, colaboração e desenvolvimento do raciocínio lógico, elas tendem a não priorizar de modo explícito as fragilidades que se manifestam em avaliações externas como o SAEB. Foi justamente a recorrência desse padrão, observada em investigações anteriores, que motivou a proposição do Ciclo de Remidiação Ativa (CRA), concebido para garantir que a transposição entre linguagem natural escrita e narrativa digital ocorra de forma consciente, estruturada e formativa, com atenção às transformações linguísticas implicadas pelo novo meio e pelo contexto comunicacional \cite{rojo2012}.

Pesquisas atuais também têm evidenciado a ampla disseminação de práticas de digital storytelling nas escolas, destacando seu potencial para promover autoria, engajamento e aprendizagem multimodal. A tese de \textcite{prata2024}, ao analisar qualitativamente experiências desenvolvidas em diferentes contextos educacionais do país, revela que narrativas digitais mobilizam de forma integrada texto, imagem, som e interação, configurando-se como práticas relevantes de letramentos contemporâneos. No entanto, a autora observa que, em muitas dessas experiências, a articulação entre escrita, estrutura narrativa e processos de formalização permanece implícita, dependendo fortemente da mediação docente. No contexto do Ciclo de Remidiação Ativa, esses achados reforçam a necessidade de uma metodologia que organize de modo intencional a articulação entre linguagens e mídias, garantindo que a multimodalidade não dilua, mas potencialize o desenvolvimento da linguagem escrita e do Pensamento Computacional.

É nesse contexto que o conceito de remidiação se torna particularmente relevante. Conforme definido por \textcite{bolter2000}, a remidiação refere-se ao processo pelo qual uma mídia é reinterpretada e reconfigurada em outra. No âmbito do Ciclo de Remidiação Ativa (CRA), essa noção é apropriada pedagogicamente para designar a transposição intencional de um texto produzido inicialmente em linguagem natural escrita, em norma padrão, para uma narrativa digital interativa. Tal transposição não se reduz a uma adaptação técnica entre suportes, mas configura-se como um processo cognitivo e autoral que exige reestruturação narrativa, escolhas semióticas e formalização progressiva das ideias.

Essa compreensão da remidiação como processo transformador fundamenta a organização do CRA como um ciclo formativo, e não como uma sequência linear de etapas. As fases que o compõem são interdependentes e admitem retornos, revisões e refinamentos ao longo do percurso. De forma sintética, o ciclo compreende: (i) a problematização do tema e do gênero textual; (ii) a produção autoral do texto em linguagem natural escrita, com atenção aos critérios de coesão, coerência e adequação linguística; (iii) a decomposição da narrativa em unidades estruturais, como cenas e ações; (iv) a construção de um algoritmo narrativo em linguagem natural; (v) a implementação da narrativa digital em ambiente de programação visual; e (vi) o compartilhamento e a socialização das produções.

Cada uma dessas etapas mobiliza diferentes dimensões do Pensamento Computacional e da criatividade, sempre ancoradas na intencionalidade comunicativa que orienta a produção narrativa. A escrita em linguagem natural, trabalhada de forma consciente e reflexiva, especialmente no que se refere ao uso da norma padrão e das variedades linguísticas, constitui um eixo estruturante do CRA, podendo atuar tanto como ponto de partida quanto como instância de retorno e refinamento ao longo do ciclo. Dessa forma, o Ciclo de Remidiação Ativa não corresponde a uma sequência didática única, mas a uma metodologia formativa que admite múltiplas sequências didáticas em seu interior, ajustadas aos diferentes contextos educacionais.

O caráter ativo do CRA reside no protagonismo atribuído aos sujeitos ao longo de todo o processo. Os participantes não apenas consomem narrativas ou executam tutoriais técnicos, mas produzem textos, tomam decisões autorais, enfrentam problemas de implementação e refletem sobre suas escolhas. A mediação docente, nesse contexto, assume o papel de orientação e acompanhamento do processo. Além disso, o CRA foi concebido para favorecer práticas colaborativas, nas quais a discussão coletiva e a socialização das ideias, em diferentes níveis de elaboração e formalização, integram o próprio desenvolvimento da linguagem e do Pensamento Computacional.

Assim, o Ciclo de Remidiação Ativa (CRA) configura-se como uma metodologia formativa não linear que integra, de modo orgânico, linguagem, computação e autoria digital, tendo a remidiação como gesto metodológico central. Originalmente desenvolvido para responder a duas demandas essenciais da educação contemporânea: a deficiência no ensino de língua portuguesa e a necessidade de desenvolvimento do Pensamento Computacional, o CRA também propicia a tomada de consciência sobre o uso efetivo das variedades linguísticas, incluindo a norma padrão. O percurso pode iniciar-se em diferentes práticas ou mídias, desde que o ciclo se efetive plenamente por meio da articulação explícita entre a linguagem natural escrita e a narrativa digital interativa, inclusive em movimento reverso. A Figura~\ref{fig:cra} sintetiza essa organização metodológica ao representar o CRA como um ciclo mediado por retornos e reencaminhamentos constantes, evidenciando que não se trata de uma sequência didática prescritiva, mas de um percurso reflexivo orientado à construção de sentido em múltiplas linguagens.

\begin{figure}[H]
\centering
\includegraphics[width=0.8\textwidth]{figura1.png}
\caption{Ciclo de Remidiação Ativa.}
\label{fig:cra}
\source{Elaboração própria.}
\end{figure}

\subsection{Scratch como ambiente de programação autoral}

O Scratch é uma linguagem de programação visual desenvolvida pelo MIT Media Lab com o objetivo de tornar a programação mais acessível, criativa e significativa, especialmente para crianças, jovens e educadores. Seu diferencial está na interface intuitiva baseada em blocos coloridos e encaixáveis, que elimina a necessidade de dominar sintaxes textuais complexas. Com isso, permite ao usuário focar na lógica, na organização narrativa e no design expressivo de seus projetos, favorecendo o desenvolvimento do pensamento computacional em articulação com competências cognitivas e autorais (Resnick et al., 2009).

Na proposta formativa do Ciclo de Remidiação Ativa (CRA), o Scratch, ou outra plataforma de programação, pode ser adotado como artefato digital na etapa de implementação das narrativas interativas. No caso deste estudo, após a produção textual e a decomposição narrativa, os professores participantes remidiaram suas fábulas no ambiente Scratch, criando histórias digitais com personagens, cenários, animações e comandos programados. Esse processo de transposição exige que os participantes articulem múltiplas linguagens (verbal, visual, sonora e de implementação algorítmica) favorecendo o desenvolvimento de formas híbridas de letramento.

O Scratch oferece um conjunto de recursos que ampliam as possibilidades de autoria criativa: bibliotecas de imagens e sons, importação de mídias externas, criação de personagens personalizados, inserção de efeitos sonoros e construção de interações por meio de eventos, condicionais e variáveis. Tais funcionalidades permitem a construção de narrativas não lineares, com múltiplos percursos e desfechos, possibilitando variações estéticas e lúdicas. Como ambiente de autoria, o Scratch estimula a experimentação, a expressão individual e a resolução criativa de problemas, conectando a lógica computacional à sensibilidade narrativa.

Além de suas qualidades pedagógicas, a plataforma apresenta vantagens operacionais relevantes em contextos escolares. Seu uso offline, por meio de aplicativo instalado, torna-se estratégico em redes com acesso limitado à internet. A existência de tutoriais em português, ampla documentação disponível e uma comunidade ativa de educadores contribuem para a autonomia dos usuários e a continuidade do uso após a formação inicial.

Dessa forma, o Scratch é utilizado como um ambiente formativo que integra criação, lógica e estética em uma mesma experiência. Atua como mediador simbólico e cognitivo, alinhado à perspectiva histórico-cultural de \textcite{vygotsky2007}, e viabiliza a concretização dos princípios dos multiletramentos e da remidiação criativa no contexto da formação docente.


\section{Metodologia}\label{sec-metodologia}

A investigação sobre a manifestação da criatividade nas produções autorais dos professores em formação, a partir da reescrita de fábulas criadas por \textcite{foohs2022b}, no contexto do Ciclo de Remidiação Ativa (CRA), adotou uma abordagem qualitativa, de natureza descritiva e interpretativa. Essa opção metodológica responde à complexidade das narrativas digitais interativas como objetos multimodais e situados, que articulam linguagem escrita, estrutura narrativa, lógica algorítmica e design computacional. Conforme \textcite{bogdan1994}, a pesquisa qualitativa permite a análise aprofundada de artefatos culturais em seus contextos naturais, considerando seus aspectos estruturais, simbólicos e formativos. Nessa direção, as produções foram examinadas como manifestações concretas de práticas criativas mediadas por tecnologia, revelando dimensões relevantes dos processos de autoria e do desenvolvimento profissional docente.

Para garantir densidade interpretativa, foi elaborada uma matriz analítica híbrida, que articula os critérios de criatividade de \textcite{torrance1995} (fluência, originalidade, elaboração e flexibilidade) com os indicadores de criatividade computacional descritos por \textcite{brennan2012}, que compreendem práticas, conceitos e perspectivas do pensamento computacional. Esses referenciais são complementados pelos parâmetros de organização algorítmica propostos pela Nota Técnica nº 21/2024 do CIEB.

A pesquisa foi aprovada pelo Comitê de Ética em Pesquisa da universidade envolvida, conforme parecer CAEE 77530824.0.0000.5347, de acordo com a Resolução nº 510/2016 do Conselho Nacional de Saúde.


\subsection{Contexto, participantes e corpus}

O estudo foi desenvolvido no primeiro semestre de 2025, no contexto de uma formação continuada promovida por uma rede pública de ensino. A formação, baseada no Ciclo de Remidiação Ativa (CRA), teve duração de três meses e foi estruturada em nove módulos (três presenciais e seis online), mediados pela plataforma Google Sala de Aula.

Participaram 48 professores da rede pública, atuantes nos anos finais do Ensino Fundamental, com adesão voluntária via edital interno. Para fins de pesquisa, foram consideradas as produções de 16 docentes que assinaram o Termo de Consentimento Livre e Esclarecido (TCLE).

Dentre esses 16 participantes, apenas um declarou nunca ter tido contato prévio com o ambiente Scratch, o que exigiu da equipe formadora uma mediação mais cuidadosa e acompanhamento diferenciado, especialmente durante as etapas de decomposição narrativa e implementação digital.

O conjunto de dados analisado neste estudo é composto pelas produções autorais entregues pelos participantes ao longo das etapas do Ciclo de Remidiação Ativa (CRA). Essas produções constituem o corpus da pesquisa e correspondem a diferentes momentos do percurso formativo, podendo variar em número conforme a participação e a completude das entregas.

As três categorias de produções são:

\begin{enumerate}[label=\alph*)]
\item	Fábula reescrita: texto autoral baseado na releitura de uma fábula.
\item	Algoritmo narrativo: decomposição da narrativa em estrutura sequencial, com foco na lógica computacional.
\item	Projeto no Scratch: implementação digital da narrativa como experiência interativa multimodal, com uso de programação visual em blocos.
\end{enumerate}

Nem todos os participantes completaram as três tarefas nos nove módulos da formação, mas todas as produções entregues foram consideradas válidas, respeitando a diversidade dos percursos formativos. Essa abordagem metodológica permite reconhecer manifestações criativas mesmo quando o Ciclo não é concluído integralmente.

\subsection{Procedimentos e matriz de análise}

As produções foram organizadas em três categorias: fábula reescrita, algoritmo narrativo e projeto no Scratch. Cada tipo de produção foi analisado segundo critérios específicos:

\begin{enumerate}[label=\alph*)]
\item	Fábulas reescritas: \textcite{torrance1995} – fluência, originalidade, elaboração, flexibilidade;
\item	Algoritmos narrativos: \textcite{cieb2024} – organização sequencial, clareza, coerência lógica, adaptabilidade;
\item	Projetos no Scratch: \textcite{brennan2012} – práticas (iterar, modularizar, depurar), conceitos (eventos, variáveis, paralelismo), perspectivas (autoria, expressividade, compartilhamento).
\end{enumerate}

A análise foi qualitativa, descritiva e interpretativa, com registros sistemáticos em matrizes analíticas baseadas nos critérios de criatividade textual e computacional, favorecendo a identificação de tendências, variações e modos expressivos de autoria nas diferentes produções. Na \Cref{quadro:matriz}, apresenta-se a matriz analítica de criatividade utilizada:

\begin{table}[ht]
\centering
\caption{Matriz de análise.}
\label{quadro:matriz}

\begin{tabular}{p{3.5cm}p{3.5cm}p{6.5cm}}
\toprule
Objeto de Análise & Referencial & Critérios Analíticos \\
\midrule
Fábula reescrita & \textcite{torrance1995} & Fluência, originalidade, elaboração, flexibilidade \\
Algoritmo narrativo & \textcite{cieb2024} & Organização sequencial, clareza, coerência lógica, adaptabilidade \\
Projeto no Scratch & \textcite{brennan2012} & Práticas (iterar, modularizar, depurar), conceitos (eventos, variáveis, paralelismo), perspectivas (autoria, expressividade, compartilhamento) \\
\bottomrule
\end{tabular}

\vspace{0.3cm}
\source{Elaboração própria.}
\end{table}

\subsection{Limites e possibilidades da pesquisa}

Entre os limites da investigação, destaca-se a análise centrada somente nas produções autorais dos participantes, realizadas durante a formação. Embora tenha sido aplicado um questionário final com questões abertas e fechadas, esse instrumento não foi utilizado como fonte de dados para este artigo, uma vez que as respostas foram anônimas e não permitiriam associação direta com os projetos analisados. Assim, a ausência de triangulação com entrevistas, observações ou instrumentos nominativos delimita o escopo interpretativo dos achados. Além disso, o corpus se apresenta heterogêneo, com entregas parciais e formatos variados, o que reflete as condições reais da formação mediada por tecnologias e os tempos disponíveis dos professores em exercício.

Apesar dessas limitações, a pesquisa oferece uma leitura aprofundada das potencialidades formativas do Ciclo de Remidiação Ativa (CRA), evidenciando sua viabilidade como metodologia replicável que articula linguagem, lógica, programação e estética em práticas criativas de retextualização. Futuros estudos poderão aprofundar esses achados por meio de estratégias complementares de coleta e análise, incluindo o cruzamento de dados com instrumentos identificados e o acompanhamento longitudinal das reverberações na prática docente.

\section{Apresentação e análise dos resultados}\label{sec-resultados}

Esta seção apresenta os resultados da análise qualitativa das produções autorais desenvolvidas pelos professores durante o Ciclo de Remidiação Ativa (CRA), com foco na manifestação da criatividade em três dimensões expressivas: (a) verbal, por meio das fábulas reescritas, analisadas com base nos critérios de \textcite{torrance1995}; (b) estrutural-algorítmica, por meio dos scripts narrativos, avaliados segundo os princípios de clareza, coerência e sequenciamento lógico propostos pelo \textcite{cieb2024}; e (c) computacional, por meio dos projetos implementados no Scratch, interpretados com base nos indicadores de \textcite{brennan2012}, que contemplam práticas, conceitos e perspectivas do pensamento computacional criativo.

A partir dessa matriz analítica, os 16 participantes foram agrupados em três níveis de manifestação criativa: excelente (6 docentes), satisfatório (7 docentes) e elementar (3 docentes). Optou-se por manter todas as produções na análise, independentemente do grau de completude, a fim de preservar a diversidade dos percursos formativos e reconhecer diferentes modos de autoria emergentes.

Cada subseção descreve os padrões observados em cada grupo, destacando recorrências e singularidades nas produções. Essa abordagem busca oferecer uma leitura contextualizada e sensível das formas pelas quais os professores mobilizaram linguagem, lógica, programação e estética em suas produções, revelando trajetórias singulares de autoria mediada por tecnologias digitais.

\subsection{Grupo 1 – Produções com desempenho excelente (n=6)}

O Grupo 1 reúne participantes cujas produções demonstraram manifestações criativas de alto nível em todas as linguagens analisadas: verbal, algorítmica e computacional. As fábulas reescritas revelam sofisticação narrativa, elaboração simbólica, fluência expressiva e uso inventivo de estruturas textuais. Os algoritmos narrativos apresentam criatividade na organização lógica e na transposição da narrativa para uma estrutura programável. Já os projetos no Scratch evidenciam autoria expressiva e uso criativo dos recursos computacionais.

As fábulas produzidas pelos participantes do Grupo 1 – Nível Excelente revelam alto grau de criatividade verbal escrita, com forte presença de fluência, originalidade, elaboração e flexibilidade. Observam-se recursos como ironia, poesia rimada, bifurcação de finais, metalinguagem e crítica social, todos articulados com coesão interna e marcas autorais nítidas.

Os algoritmos narrativos desenvolvidos por esse grupo demonstram criatividade na estruturação sequencial e na representação funcional da narrativa. Há uso expressivo de listas, ramificações condicionais, sorteios, interações e comandos diversos, que conferem dinamicidade à lógica narrativa e apontam para soluções inovadoras na transposição para o ambiente computacional.

Nesse sentido, as implementações computacionais desenvolvidas por esse grupo evidenciam criatividade na estruturação das interações e na materialização da lógica narrativa no ambiente Scratch. Os projetos revelam escolhas autorais marcantes, organização estética, experimentação narrativa e uso criativo dos recursos computacionais (eventos, variáveis, paralelismo, laços, entre outros). A perspectiva computacional é mobilizada de modo expressivo e intencional, resultando em experiências interativas originais e coerentes.


\subsection{Grupo 2 – Produções com desempenho satisfatório (n=7)}

O Grupo 2 é composto por participantes cujas produções apresentaram manifestações criativas com oscilações na fluência narrativa, na elaboração estética e na exploração dos recursos disponíveis em cada linguagem. As criações revelam intenção autoral e tentativas de inovação, ainda que com limitações estruturais, simbólicas ou técnicas. Há recorrência de temas como humor, acessibilidade, crítica e metáfora, porém com menor densidade narrativa e estilística em relação ao Grupo 1 – Nível Excelente. 

De modo geral, o Grupo 2 – Nível Satisfatório representa um nível intermediário de expressão criativa, com pistas de autoria e exploração estética, apesar das restrições técnicas e narrativas. As produções revelam criatividade em processo, com caminhos diversos de apropriação, experimentação e desenvolvimento ao longo da formação.

\subsection{Grupo 3 – Produções com desempenho elementar (n=3)}

O Grupo 3 reúne participantes cujas produções demonstraram manifestações criativas bastante incipientes, com baixo grau de elaboração narrativa, expressividade estética limitada e uso muito restrito dos recursos computacionais. As entregas apresentam lacunas significativas nas três linguagens analisadas (verbal, algorítmica e computacional), com predomínio de fragmentos descritivos, estruturas incompletas ou desarticuladas e ausência de traços autorais consistentes. Ainda assim, alguns registros pontuais indicam sinais embrionários de intenção narrativa, que serão observados nos comentários analíticos a seguir.

O Grupo 3 – Nível Elementar apresenta tentativas de construção de algoritmos narrativos, revelando um nível elementar de apropriação da lógica narrativa e dos princípios da decomposição algorítmica. As entregas consistem em descrições genéricas ou frases desconexas, sem encadeamento funcional entre as ações e falas, o que inviabiliza a transposição para ambientes de programação como o Scratch. A ausência de algoritmos viáveis compromete diretamente a articulação entre linguagem verbal e computacional, interrompendo o Ciclo de Remidiação Ativa e revelando dificuldades conceituais na organização lógica da narrativa.

Os projetos computacionais analisados revelam implementações rudimentares, com baixo nível de interatividade, ausência de modularização e pouco ou nenhum uso dos conceitos fundamentais da programação por blocos. Em geral, os personagens não executam ações significativas, não há diálogos consistentes ou sequência de eventos, e os recursos visuais e sonoros são aplicados de forma limitada. A expressividade autoral não se consolida, e a linguagem computacional não é mobilizada como instrumento de autoria narrativa.

Em síntese, as produções do Grupo 3 – Nível Elementar refletem grandes desafios na mobilização da criatividade nas três linguagens analisadas. Os dados sugerem a necessidade de intensificação do acompanhamento pedagógico, ampliação do tempo de elaboração e fortalecimento da mediação formativa, para que esses sujeitos possam desenvolver repertório autoral, narrativo e computacional com maior consistência e intencionalidade.

\subsection{Síntese interpretativa e cruzamento dos dados}

A análise comparativa das produções digitais, organizadas em três grupos de desempenho, evidenciou padrões distintos de manifestação da criatividade docente, articulando linguagem escrita, abstração algorítmica e programação interativa.

No Grupo 1 – Nível Excelente, destacaram-se os quatro critérios de Torrance: fluência, originalidade, elaboração simbólica e flexibilidade. As fábulas apresentaram metáforas inventivas e personagens densos; os algoritmos, estrutura lógica com bifurcações; e os projetos no Scratch, domínio técnico e autoralidade estética, com uso expressivo da linguagem de blocos.

O Grupo 2 – Nível Satisfatório revelou tentativas de inovação (humor, crítica, lirismo), ainda que com menor densidade simbólica. Os algoritmos demonstraram lógica básica e os projetos computacionais, apropriação inicial dos conceitos, embora com limitações técnicas e autorais. Observa-se esforço criativo em diferentes níveis.

Já o Grupo 3 – Nível Elementar apresentou manifestações criativas incipientes. As fábulas tendem à reprodução esquemática; os algoritmos são ausentes ou rudimentares; e os projetos no Scratch carecem de interatividade e simbolismo. Ainda assim, alguns registros indicam intenção narrativa embrionária.

Transversalmente, percebe-se que a criatividade não depende apenas do domínio técnico, mas emerge da interação entre linguagem, imaginação e lógica narrativa. A estrutura formativa, com transposições sucessivas entre linguagens (natural, algorítmica, computacional), favoreceu pensamento divergente, exploração simbólica e autoria multimodal. A reelaboração contínua expandiu o repertório expressivo dos participantes e tensionou positivamente suas zonas de desenvolvimento. Essa dinâmica corrobora a perspectiva histórico-cultural de \textcite{vygotsky2007}, ao evidenciar que a aprendizagem mediada por instrumentos semióticos, neste caso, a linguagem escrita, o algoritmo narrativo e a programação, impulsiona o desenvolvimento, ao mesmo tempo em que se articula a práticas de letramentos múltiplos, conforme proposto por \textcite{rojo2012}.

Os referenciais analíticos de \textcite{torrance1995,brennan2012,cieb2024} permitiram identificar práticas autorais marcadas por inovação temática, sofisticação narrativa, expressividade digital e intencionalidade estética, mesmo quando incipientes, como sinais de criatividade em processo.

Esses achados permitem compreender o Ciclo de Remidiação Ativa (CRA) como um gesto metodológico estruturante que explica as diferenças observadas entre os grupos analisados. Mais do que um conjunto de etapas, o CRA opera como um dispositivo formativo que organiza a passagem consciente entre linguagem natural escrita, abstração algorítmica e implementação computacional, criando condições para a emergência da criatividade como processo autoral. Nos casos em que o ciclo se efetivou de modo mais consistente, observou-se maior densidade narrativa, sofisticação simbólica e integração entre lógica e expressividade; quando o ciclo foi interrompido ou fragilizado, as produções tenderam à reprodução esquemática e à limitação estrutural.

A articulação dos referenciais teóricos adotados mostrou-se fundamental para a leitura integrada dos dados. Os critérios de \textcite{torrance1995} permitiram identificar manifestações de fluência, originalidade, elaboração e flexibilidade na linguagem verbal; os indicadores do \textcite{cieb2024} orientaram a análise da organização lógica e da decomposição narrativa; e o modelo de \textcite{brennan2012} possibilitou compreender como práticas, conceitos e perspectivas computacionais se manifestaram no Scratch como linguagem de autoria. Esses referenciais não operaram de forma isolada, mas de maneira complementar, sendo integrados pelo movimento de remidiação que caracteriza o CRA.

Assim, os resultados indicam que a criatividade docente observada não decorre prioritariamente do domínio técnico da programação em blocos, mas da possibilidade de reconfigurar ideias, sentidos e narrativas ao longo de sucessivas transposições entre linguagens. Nesse processo, o trabalho consciente com a linguagem escrita mostrou-se decisivo, não como formalismo normativo, mas como prática reflexiva que sustenta a reorganização de sentidos, em consonância com a pedagogia da variação linguística proposta por \textcite{faraco2015}.

O CRA evidencia-se, portanto, como uma metodologia formativa potencialmente capaz de articular criatividade e Pensamento Computacional de modo situado, reflexivo e não reducionista, evitando tanto a diluição conceitual do PC quanto sua redução a habilidades instrumentais (\Cref{tab:sintese}).

\begin{table}[ht]
\centering
\caption{Síntese da manifestação criativa por Grupo e Linguagem.}
\label{tab:sintese}

\small
\setlength{\tabcolsep}{3pt}

\resizebox{\textwidth}{!}{%
\begin{tabular}{p{2.2cm}p{3cm}p{3cm}p{3.4cm}p{2.8cm}}
\toprule
Grupo (N=16) & Linguagem Verbal (Torrance) & Linguagem Algorítmica (CIEB) & Linguagem Computacional (B\&R) & Implicação \\
\midrule
Excelente (n=6) & Alta fluência e flexibilidade & Lógica clara. Uso de ramificações & Alto domínio de conceitos (variáveis, paralelismo) e práticas (modularização) & Transposição criativa integrada \\
Satisfatório (n=7) & Elaboração oscilante. Originalidade temática (humor, inclusão) & Lógica básica. Estruturas genéricas. Falta de precisão algorítmica & Uso funcional de eventos e falas; ausência de paralelismo ou estruturas complexas & Criatividade limitada pela estrutura lógica \\
Elementar (n=3) & Reprodução esquemática. Baixa fluência & Ausente ou rudimentar (interrupção do CRA) & Não-interativo, rudimentar e ausência de conceitos-chave & Falha na abstração lógica e na construção criativa \\
\bottomrule
\end{tabular}%
}

\vspace{0.3cm}
\source{Elaboração própria.}
\end{table}

Por fim, os resultados convergem com os objetivos da BNCC \cite{bncc2018,cne2022}, especialmente os códigos EF69LP18/19 (Linguagem) e EF69CO03/05 (Computação), apontando o potencial de propostas interdisciplinares para fomentar a autoria criativa na formação docente.

\section{Considerações finais}\label{sec-conclusoes}

Este artigo analisou a criatividade docente na articulação entre linguagem, narrativa e Pensamento Computacional a partir do Ciclo de Remidiação Ativa (CRA), tomando como corpus produções autorais desenvolvidas em formação continuada com narrativas digitais interativas no Scratch. A opção por abordagem qualitativa interpretativa fundamentou-se na compreensão de que criatividade, autoria e processos de formalização narrativa não se expressam de modo mensurável por indicadores isolados, mas emergem de trajetórias formativas, escolhas autorais e transposições sucessivas entre linguagens.

A análise das produções evidenciou progressões formativas distintas nas dimensões verbal, algorítmica e computacional, organizadas em três níveis de manifestação criativa. Nos casos em que o CRA se efetivou de maneira consistente, observaram-se maior densidade narrativa, fluência expressiva, flexibilidade criativa e integração entre lógica e expressividade computacional. Já nas produções mais incipientes, a interrupção ou fragilização do Ciclo revelou dificuldades na abstração, na decomposição narrativa e na formalização das ideias, comprometendo a autoria digital. Esses resultados reforçam a compreensão de que a criatividade não decorre do domínio técnico da programação, mas da possibilidade de reorganizar sentidos ao longo de sucessivas remidiações.

O cruzamento entre os referenciais de \textcite{torrance1995,brennan2012,cieb2024} mostrou-se fecundo para interpretar os dados de forma integrada. Enquanto os critérios de Torrance permitiram identificar manifestações de criatividade na linguagem verbal, os indicadores do CIEB orientaram a análise da organização lógica e da clareza estrutural, e o modelo de Brennan e Resnick possibilitou compreender o Scratch como espaço de autoria computacional, no qual práticas, conceitos e perspectivas se articulam. O CRA operou como eixo integrador desses referenciais, funcionando não como sequência didática prescritiva, mas como gesto metodológico que organiza a passagem consciente entre linguagens.

Nesse sentido, a análise realizada não buscou classificar sujeitos ou validar instrumentos, mas compreender processos formativos em curso. Os quadros analíticos utilizados ao longo do estudo tiveram a função de evidenciar padrões e recorrências, e não de esgotar a interpretação dos dados. A síntese interpretativa permitiu articular essas evidências à fundamentação teórica, indicando que a criatividade docente se constrói como prática mediada, situada e culturalmente orientada, em consonância com uma perspectiva histórico-cultural de aprendizagem.

Como limitação, reconhece-se que os dados analisados derivam exclusivamente das produções entregues, sem o apoio de entrevistas, diários reflexivos ou outros instrumentos que permitissem acessar percepções subjetivas dos participantes sobre o processo formativo. Pesquisas futuras poderão integrar essas dimensões, bem como aplicar o CRA em contextos escolares diversos, com diferentes públicos e tempos de intervenção, aprofundando a análise de suas potencialidades e desafios.

Conclui-se que a integração entre linguagem, narrativa e computação, quando orientada por uma metodologia formativa baseada na remidiação consciente entre linguagens, constitui um caminho promissor para a formação docente. O Ciclo de Remidiação Ativa demonstra potencial para enfrentar simultaneamente fragilidades no uso da língua portuguesa e demandas contemporâneas relacionadas ao desenvolvimento do Pensamento Computacional, sem reduzir nenhum desses campos a abordagens instrumentais. Ao valorizar autoria, reflexão e mediação pedagógica, o CRA afirma-se como proposta metodológica consistente para práticas formativas críticas, criativas e culturalmente situadas.

\printbibliography\label{sec-bib}
% if the text is not in Portuguese, it might be necessary to use the code below instead to print the correct ABNT abbreviations [s.n.], [s.l.]
%\begin{portuguese}
%\printbibliography[title={Bibliography}]
%\end{portuguese}


%full list: conceptualization,datacuration,formalanalysis,funding,investigation,methodology,projadm,resources,software,supervision,validation,visualization,writing,review
\begin{contributors}[sec-contributors]
\authorcontribution{Marcelo Magalhães Foohs}[conceptualization,datacuration,formalanalysis,methodology,projadm,writing]
\authorcontribution{Priscila Castioni Isele}[conceptualization,datacuration,review]
\authorcontribution{Randerson Oliveira Melville Rebouças}[conceptualization,datacuration,review]
\end{contributors}

\begin{dataavailability}
\txtdataavailability{dataonly} % options: dataavailable, dataonly, databody, datanotav, nodata
\end{dataavailability}

\end{document}
